% GENERATED FILE. Edit canonical lessons, then run npm run generate:books.
% canonical-source-hash: fnv1a64:b17d3f88822760d6
% canonical-lessons: GE-C464-hitzewelle, GE-C464-duerre, GE-C464-luftfeuchtigkeit, GE-C464-schwuel, GE-C464-austrocknen

\chapter{Die Hitzewelle}
\label{ch:ge-die-hitzewelle}
\hlchaptermodality{464}

\begin{chapteropening}
\textbf{By the end of this chapter:} \emph{I can talk about hot, dry weather: the heatwave, the drought, the humidity, muggy air and drying out.}

Complete the last lesson of chapter 464: I can talk about hot, dry weather: the heatwave, the drought, the humidity, muggy air and drying out.
\end{chapteropening}

\section[die Hitzewelle]{die Hitzewelle — the heatwave}

\label{lesson:GE-C464-hitzewelle}



\begin{quote}
Before the new one: say the German for the thaw, then the German for to shiver.
\end{quote}

\subsection*{You'll want to know: die Hitzewelle}
\textbf{die Hitzewelle} — \textquotedblleft{}the heatwave\textquotedblright{}.

\begin{quote}
Während der Hitzewelle sind die Schwimmbäder voll.
\end{quote}

\begin{quote}
\emph{During the heatwave the swimming pools are full.}
\end{quote}

Plural: \emph{die Hitzewellen}.

\begin{cousinweb}
\emph{Hitzewelle} is built from \emph{Hitze} 'heat', an inherited word related to \emph{heiß} 'hot', and \emph{Welle} 'wave', also an inherited Germanic word (EtymDB lists \emph{Welle}).
\end{cousinweb}

\begin{grammarlens}[title={What you've built}]
One more word for telling, explaining and describing at B1.
\end{grammarlens}

\subsection*{Your turn}
\begin{itemize}
\raggedright
  \item \emph{Say it:} \emph{die Hitzewelle}
  \item \emph{Say it:} \emph{die Hitzewelle}, once more
  \item \emph{Say it:} a sentence of your own with \emph{die Hitzewelle}
  \item \emph{Recall:} say the German for the thaw, then the German for to shiver, then say \emph{die Hitzewelle} again
\end{itemize}

\begin{tcolorbox}[breakable,colback=teal!4,colframe=teal!35!black,title={Before you move on}]
What does die Hitzewelle mean? (\textquotedblleft{}The heatwave\textquotedblright{}.) Say it once more.
\end{tcolorbox}
\section[die Dürre]{die Dürre — the drought}

\label{lesson:GE-C464-duerre}



\begin{quote}
Before the new one: say the German for to shiver, then the German for the heatwave.
\end{quote}

\subsection*{You'll want to know: die Dürre}
\textbf{die Dürre} — \textquotedblleft{}the drought\textquotedblright{}.

\begin{quote}
Nach der langen Dürre ist das Gras braun.
\end{quote}

\begin{quote}
\emph{After the long drought the grass is brown.}
\end{quote}

Plural: \emph{die Dürren}. \emph{Dürr} — very dry; thin.

\begin{cousinweb}
\emph{Dürre} is built on \emph{dürr} 'dry, withered', an inherited Germanic adjective related to \emph{Durst} 'thirst'. EtymDB does not list the noun.
\end{cousinweb}

\begin{grammarlens}[title={What you've built}]
One more word for telling, explaining and describing at B1.
\end{grammarlens}

\subsection*{Your turn}
\begin{itemize}
\raggedright
  \item \emph{Say it:} \emph{die Dürre}
  \item \emph{Say it:} \emph{die Dürre}, once more
  \item \emph{Say it:} a sentence of your own with \emph{die Dürre}
  \item \emph{Recall:} say the German for to shiver, then the German for the heatwave, then say \emph{die Dürre} again
\end{itemize}

\begin{tcolorbox}[breakable,colback=teal!4,colframe=teal!35!black,title={Before you move on}]
What does die Dürre mean? (\textquotedblleft{}The drought\textquotedblright{}.) Say it once more.
\end{tcolorbox}
\section[die Luftfeuchtigkeit]{die Luftfeuchtigkeit — the humidity}

\label{lesson:GE-C464-luftfeuchtigkeit}



\begin{quote}
Before the new one: say the German for the heatwave, then the German for the drought.
\end{quote}

\subsection*{You'll want to know: die Luftfeuchtigkeit}
\textbf{die Luftfeuchtigkeit} — \textquotedblleft{}the humidity\textquotedblright{}.

\begin{quote}
Die Luftfeuchtigkeit ist heute sehr hoch.
\end{quote}

\begin{quote}
\emph{The humidity is very high today.}
\end{quote}

No plural.

\begin{cousinweb}
\emph{Luftfeuchtigkeit} is built from \emph{Luft} 'air', inherited from Proto-Germanic \emph{luftuz} (EtymDB), and \emph{Feuchtigkeit} 'dampness', from \emph{feucht} 'damp'.
\end{cousinweb}

\begin{grammarlens}[title={What you've built}]
One more word for telling, explaining and describing at B1.
\end{grammarlens}

\subsection*{Your turn}
\begin{itemize}
\raggedright
  \item \emph{Say it:} \emph{die Luftfeuchtigkeit}
  \item \emph{Say it:} \emph{die Luftfeuchtigkeit}, once more
  \item \emph{Say it:} a sentence of your own with \emph{die Luftfeuchtigkeit}
  \item \emph{Recall:} say the German for the heatwave, then the German for the drought, then say \emph{die Luftfeuchtigkeit} again
\end{itemize}

\begin{tcolorbox}[breakable,colback=teal!4,colframe=teal!35!black,title={Before you move on}]
What does die Luftfeuchtigkeit mean? (\textquotedblleft{}The humidity\textquotedblright{}.) Say it once more.
\end{tcolorbox}
\section[schwül]{schwül — muggy, sultry}

\label{lesson:GE-C464-schwuel}



\begin{quote}
Before the new one: say the German for the drought, then the German for the humidity.
\end{quote}

\subsection*{You'll want to know: schwül}
\textbf{schwül} — \textquotedblleft{}muggy, sultry\textquotedblright{}.

\begin{quote}
Vor dem Gewitter ist es oft schwül.
\end{quote}

\begin{quote}
\emph{Before a thunderstorm it is often muggy.}
\end{quote}

\emph{Die Schwüle} — the muggy heat.

\begin{cousinweb}
\emph{schwül} came into the standard language from Low German; it is usually linked to \emph{schwelen} 'to smoulder'. EtymDB lists it with the sense 'sultry' but gives no source.
\end{cousinweb}

\begin{grammarlens}[title={What you've built}]
One more word for telling, explaining and describing at B1.
\end{grammarlens}

\subsection*{Your turn}
\begin{itemize}
\raggedright
  \item \emph{Say it:} \emph{schwül}
  \item \emph{Say it:} \emph{schwül}, once more
  \item \emph{Say it:} a sentence of your own with \emph{schwül}
  \item \emph{Recall:} say the German for the drought, then the German for the humidity, then say \emph{schwül} again
\end{itemize}

\begin{tcolorbox}[breakable,colback=teal!4,colframe=teal!35!black,title={Before you move on}]
What does schwül mean? (\textquotedblleft{}Muggy, sultry\textquotedblright{}.) Say it once more.
\end{tcolorbox}
\section[austrocknen]{austrocknen — to dry out, to dry up}

\label{lesson:GE-C464-austrocknen}



\begin{quote}
Before the new one: say the German for the humidity, then the German for muggy.
\end{quote}

\subsection*{You'll want to know: austrocknen}
\textbf{austrocknen} — \textquotedblleft{}to dry out, to dry up\textquotedblright{}.

\begin{quote}
Im August trocknet der kleine Bach fast ganz aus.
\end{quote}

\begin{quote}
\emph{In August the little stream almost dries up completely.}
\end{quote}

Separable; Perfekt with \emph{sein} when it means drying up: \emph{der Bach ist ausgetrocknet}.

\begin{cousinweb}
\emph{austrocknen} is built from \emph{aus} 'out' and \emph{trocknen} 'to dry', from \emph{trocken} 'dry', inherited from Middle High German \emph{trucken} (EtymDB).
\end{cousinweb}

\begin{grammarlens}[title={What you've built}]
One more word for telling, explaining and describing at B1.
\end{grammarlens}

\subsection*{Your turn}
\begin{itemize}
\raggedright
  \item \emph{Say it:} \emph{austrocknen}
  \item \emph{Say it:} \emph{austrocknen}, once more
  \item \emph{Say it:} a sentence of your own with \emph{austrocknen}
  \item \emph{Recall:} say the German for the humidity, then the German for muggy, then say \emph{austrocknen} again
\end{itemize}

\begin{tcolorbox}[breakable,colback=teal!4,colframe=teal!35!black,title={Before you move on}]
What does austrocknen mean? (\textquotedblleft{}To dry out, to dry up\textquotedblright{}.) Say it once more.
\end{tcolorbox}
